\section{Experiments and Results}
\label{sec:experiments}

% introduction of task being studied and description of the demonstration dataset
The proposed methods for choosing an action is evaluated for the task of tying an overhand knot using the Willow Garage PR2 robot platform.
The dataset of demonstrations that we use is the one collected by \citet{Schulmanetal_ISRR2013} for their rope-tying experiments.
The demonstrations were collected in the real world and they consist of 148 pairs of point clouds and gripper trajectories.
The point clouds of the ropes were collected using an RGBD camera and then filtered based on color to remove the background.
The gripper trajectories were recorded kinesthetically, by moving the robot's gripper to tie a knot.
The dataset contains demonstrations to tie an overhand knot and also demonstrations to recover from rope states that could lead to failure cases.

The training, validation and evaluation of the trained policy are ran in a simulation environment containing the rope and the PR2.
The rope is simulated as a linked chain of capsules with bending and torsional constraints using Bullet Physics.
To evaluate, we run the trained policy in simulation on various initial rope configurations, and measure the success rate, where we define success as tying an overhand knot in a sequence of 5 or less actions.
These initial rope states are perturbed configurations of randomly chosen initial rope states from the demonstrations.
The process of perturbing a rope state consists on taking five points uniformly spaced along the rope and dragging them uniformly at random within a radius of 15 cm in a random direction.
Notice that this is the distribution of problems that we consider in our experiments.
The randomly generated initial rope configurations mentioned hereafter are drawn from this distribution.

\subsection{Expert-labelled examples}
The expert-labelled examples consist of point clouds labelled with the action that the expert human chooses.
The point clouds in the examples are from rope states that are drawn from the distribution of initial rope configurations and the rope states resulting from simulating the chosen actions.
In particular, we simulate sequences of actions for various random initial rope configurations, with a human in the loop selecting the actions to achieve a knot-tie.
In practice, not all the actions are applied and simulated to the current state for the expert to choose from.
Instead, the actions are ranked in increasing order of registration cost with respect to the current state, under the assumption that actions with smaller registration cost are likely to be better.

We evaluate the success rate of the proposed method under the following policies: greedy, one-step lookahead with width 10, and two-step lookahead with width 5.
For this evaluation, 1000 expert-labelled examples were used.
For each of these policies, the optimization hyperparameters $C$, $D$ and $F$ are first tuned via holdout validation, using a validation set of 100 random initial rope states.
The evaluation set consists of 500 random initial rope states.
We also evaluate using the nearest neighbor policy presented in \citet{Schulmanetal_ISRR2013} as a baseline comparison.
Their policy chooses the action that has the smallest bidirectional registration cost with respect to the current state.
The success rates obtained under these policies are summarized in Table~\ref{table:performance}. Notice that our best results surpasses the baseline by XX\%.

\begin{table}
  \centering
  \begin{tabular}{lc}
    \toprule
      Policy & Test Accuracy\\
    \midrule
      Nearest neighbor \cite{Schulmanetal_ISRR2013} & XX\% \\
    \midrule
      Greedy & XX\% \\
      Lookahead (depth 1, width 10) & XX\% \\
      Lookahead (depth 2, width 5) & XX\% \\
    \bottomrule
  \end{tabular}
  \caption{Success rate of tying an overhand knot under different policies.}
  \label{table:performance}
\end{table}

The learned state values for the 500 evaluated runs for the three policies are plotted in Figure~\ref{fig:values} as a function of time steps 1 through 5.
\al{shift plot from 1 to 5.}
Notice that the state values are generally increasing as we would expect since our formulation is encoding that states closer to the goal have higher values.
Also, recall that the state value of the tied knot is close to zero. \al{cite bellman constraints equation.}
Thus, we notice that in most of the runs, the policy is able to reach the goal state in 3 steps, which is the optimal number of actions to reach the goal since the expert had to take a minimun of 3 actions to tie a knot.

\begin{figure}[h!]
  \centering
    \includegraphics[width=0.9\linewidth]{figures/placeholder.png}
  \caption{State values as a function of time steps}
  \label{fig:values}
\end{figure}

We also study the relation between the number of expert-labelled examples and the success rate for the four policies.
These success rates are shown in Figure~\ref{fig:number_examples}.
The success rate of the baseline policy is a constant since the nearest neighbor policy doesn't use the expert-labelled examples.
As expected, the success rate increases as more labelled-examples are used.

\begin{figure}[h!]
  \centering
    \includegraphics[width=0.9\linewidth]{figures/placeholder.png}
  \caption{Success rates as a function of the number of expert-labelled examples for the four policies}
  \label{fig:number_examples}
\end{figure}

\subsection{Leave-One-Out Labeling}

We now evaluate the performance of training using leave-one-out labeled
examples, as outlined in Section~\ref{subsec:lool}.

After holdout validation and evaluation, done in the same manner as our man
experiment, we find that using a leave-one-out labeled example set achieves a
performance of XX\%, a drop of XX\% from our supervised result. \et{It feels
  difficult to do further analysis without an actual number here. I'll save this
  for later, if that's alright with everyone.}
